% TOP_PROBLEM: {"id": 30006674, "title": "Navier-Stokes Existence and Smoothness: Fefferman A", "category_id": 9, "category": {"id": 9, "name": "pde", "display_name": "Partial Differential Equations", "description": "PDEs and their applications in physics and geometry.", "slug": "pde", "order_index": 9, "created_at": "2026-07-31T15:26:25.671Z"}, "status": "open"}
% FIRST_POSTED: 2026-09-27
% ATTEMPT_STATUS: unresolved
% !TeX program = lualatex
\documentclass[11pt]{article}
\usepackage[margin=25mm]{geometry}
\usepackage{fontspec}
\setmainfont{Latin Modern Roman}
\usepackage{amsmath,amssymb,mathtools}
\usepackage{unicode-math}
\setmathfont{Latin Modern Math}
\usepackage{xurl,hyperref}
\hypersetup{colorlinks=true,urlcolor=blue}
\setcounter{secnumdepth}{0}
\setlength{\parindent}{0pt}
\setlength{\parskip}{5pt}
\setlength{\emergencystretch}{3em}
\begin{document}
\section*{\texorpdfstring{Navier-Stokes Existence and Smoothness: Fefferman A}{Navier-Stokes Existence and Smoothness: Fefferman A}}\label{30006674}
\subsection{Definitions and mathematical statement}
A vector field $u_0\in\mathcal S({\mathbb{R}}^3;{\mathbb{R}}^3)$ is Schwartz if $\sup_x(1+|x|)^N|\partial^\alpha u_0(x)|<\infty$ for every multi-index $\alpha$ and integer $N\ge0$. For every such field with $\nabla\cdot u_0=0$ and every $\nu>0$, seek real smooth $u,p$ on ${\mathbb{R}}^3\times[0,\infty)$ satisfying
\[
 \partial_tu+(u\cdot\nabla)u=\nu\Delta u-\nabla p,\qquad
 \nabla\cdot u=0,\qquad u(\cdot,0)=u_0,
\]
\[
 \sup_{t\ge0}\int_{{\mathbb{R}}^3}|u(x,t)|^2\,dx<\infty.
\]
Smoothness is required at time zero as well. No forcing, periodic domain, or smallness hypothesis is allowed in this selected statement, Fefferman (A).
\par\smallskip\textit{Formulation and concept references:} \href{https://www.claymath.org/wp-content/uploads/2022/06/navierstokes.pdf}{[S1]}.

\subsection{Short English statement}
Must every smooth, rapidly decreasing initial flow in unforced three-dimensional viscous fluid motion remain smooth forever with bounded total energy?
\subsection{Research attempt}

\paragraph{Scope and current literature check (27 September 2026).}
The quantifiers are: for every real divergence-free Schwartz datum on
 $\mathbb R^3$ and every fixed $\nu>0$, construct a smooth global solution
 of the \emph{unforced} equation with bounded kinetic energy. The selected
 statement is Fefferman (A), not the assertion that some allowed alternative
 in the Millennium formulation has been established.

This distinction is essential for the supplied recent sources. Theorem 1.1
 on page 1 of the analytical paper [S4] states a construction with zero
 initial datum and a smooth compactly supported external force, leading to
 blowup and alternatives (C) and (D). The announcement [S2] describes that
 forced result. Clay's 11 September 2026 announcement [S3] discusses the
 announced advance and its evaluation process. Reading these statements
 does not verify the complete proof, and their forced hypotheses do not
 supply a counterexample to the unforced target here. No associated Lean
 files were executed or treated as independent certification.

For the present argument, local strong-solution theory is the input supplied
 by Kato [E1]. The endpoint theorem of Escauriaza--Seregin--\v Sver\'ak [E2]
 says, in particular, that a suitable finite-energy Cauchy solution with
 bounded $L^\infty_tL^3_x$ norm cannot have a finite-time singularity.
 This is a conditional theorem, not an available bound for arbitrary data.
 Tao's averaged-equation construction [E3] warns that energy cancellation
 alone does not distinguish the actual Navier--Stokes nonlinearity from
 every blowup model. The calculations below reproduce classical estimates
 and test their exact failure point; no novelty is claimed for the
 small-data or continuation criteria.

\paragraph{Attack map and conventions.}
The proof route is to control vortex stretching using energy and enstrophy,
 then test whether a scale-invariant small-data argument extends to large
 data. The counterexample route is to examine whether a forced singular
 trajectory can be converted into an unforced one by removing its force.
 A third, conditional route isolates the accumulated positive strain that
 would have to be controlled by a sharper geometric argument.

Let $[0,T_*)$ be the maximal smooth strong-solution interval issued from
 $u_0$. All identities below are used before $T_*$, where Sobolev regularity
 and spatial cutoffs justify integration by parts. Write
\[
 \omega=\nabla\times u,\quad
 S=\tfrac12(\nabla u+\nabla u^{\mathsf T}),\quad
 E=\|u\|_2^2,\quad Y=\|\omega\|_2^2,\quad
 Z=\|\nabla\omega\|_2^2,\quad N=\int_{\mathbb R^3}\omega\cdot S\omega\,dx.
\]
For divergence-free fields, Plancherel gives
 $Y=\|\nabla u\|_2^2$ and $Z=\|\Delta u\|_2^2$.
 The pressure is recovered, up to a function of time, from
 $-\Delta p=\partial_i\partial_j(u_i u_j)$, with repeated indices summed.
 Pressure is not assumed to retain Schwartz decay; the target only requires
 smoothness and the velocity energy bound.

\paragraph{Exact identities and the first obstruction.}
Taking the $L^2$ product with $u$ eliminates pressure and transport because
 $\nabla\cdot u=0$. Curling the equation gives
\[
 \partial_t\omega+(u\cdot\nabla)\omega
 =\nu\Delta\omega+(\omega\cdot\nabla)u.
\]
The antisymmetric part of $\nabla u$ contributes zero to its quadratic form
 on $\omega$. Hence
\begin{equation}\label{ns:balances}
 E'=-2\nu Y,\qquad \tfrac12Y'+\nu Z=N.
\end{equation}
In particular, $\int_0^tY(s)\,ds\leq E(0)/(2\nu)$, but stretching has no
 analogous sign. Tracelessness of $S$ permits a positive eigenvalue; it
 does not imply $\omega\cdot S\omega\leq0$.

The Biot--Savart singular-integral estimate gives
 $\|\nabla u\|_3\leq C\|\omega\|_3$. Interpolation between $L^2$ and the
 three-dimensional Sobolev bound $\|\omega\|_6\leq C\|\nabla\omega\|_2$
 then gives, with one fixed universal constant $C_0$,
\begin{equation}\label{ns:stretch}
 |N|\leq\|\nabla u\|_3\|\omega\|_3^2
 \leq C_0Y^{3/4}Z^{3/4}
 \leq\tfrac\nu2 Z+C_1\nu^{-3}Y^3.
\end{equation}
The last step is Young's inequality with exponents $4/3$ and $4$.
 Thus $Y'\leq K\nu^{-3}Y^3$ for a universal $K>0$.
 This estimate gives a positive lifespan, not a bound valid for all times.

If $T_*<\infty$, bounded $Y$ would imply bounded $L^6$ velocity and allow
 continuation, as justified below. There is therefore a sequence
 $s\uparrow T_*$ with $Y(s)\to\infty$. On an interval where $Y>0$,
\[
 \frac d{dt}Y^{-2}\geq-2K\nu^{-3}.
\]
Integration from $t$ to those $s$ yields the necessary lower bound
\begin{equation}\label{ns:rate}
 Y(t)\geq\frac{\nu^{3/2}}{\sqrt{2K(T_*-t)}}.
\end{equation}
This is compatible with $\int_0^{T_*}Y<\infty$. There is no contradiction
 between the known energy budget and the derived blowup-rate condition.

To stress-test that conclusion explicitly, put
 $Y_*(t)=A(T-t)^{-1/2}$ and
\[
 E_*(t)=E_0-4\nu A\bigl(\sqrt T-\sqrt{T-t}\bigr).
\]
If $A^2\geq\nu^3/(2K)$ and $E_0>4\nu A\sqrt T$, these positive scalar
 functions satisfy $E_*'=-2\nu Y_*$ and
 $Y_*'\leq K\nu^{-3}Y_*^3$, yet $Y_*$ diverges at $T$. These are not
 velocity fields or solutions. They disprove only the proposed inference
 from these two scalar inequalities to global regularity.

\paragraph{A proved small-data subcase with the correct scaling.}
Fourier-space Cauchy--Schwarz yields
\[
 Y=\|\nabla u\|_2^2\leq\|u\|_2\|\Delta u\|_2=E^{1/2}Z^{1/2}.
\]
For $Z>0$ this implies $Y^3/Z\leq EY$. Returning to the estimate before
 Young's inequality gives the stronger, scale-sensitive statement
\begin{equation}\label{ns:invariant}
 |N|\leq C_0(EY)^{1/4}Z,
 \qquad
 \tfrac12Y'+\left(\nu-C_0(EY)^{1/4}\right)Z\leq0.
\end{equation}
If $Z=0$, finite energy makes $u$ spatially constant and hence zero, so
 there is no exceptional nonzero case.

\textbf{Proposition 5.1.}
If $C_0(E(0)Y(0))^{1/4}<\nu$, then the datum in the selected target has a
 global smooth solution, $E(t)\leq E(0)$, and $Y(t)\leq Y(0)$.

\emph{Proof.}
The zero datum gives $u=0$. Otherwise put
 $\theta_0=C_0(E(0)Y(0))^{1/4}/\nu<1$ and choose
 $\theta_1\in(\theta_0,1)$. Before the first time the coefficient ratio
 $C_0(EY)^{1/4}/\nu$ reaches $\theta_1$, equations
 \eqref{ns:balances} and \eqref{ns:invariant} make both $E$ and $Y$
 nonincreasing. Their product cannot reach that larger threshold.
 By continuity the crossing never occurs. In fact
\[
 Y(t)+2(1-\theta_0)\nu\int_0^t Z(s)\,ds\leq Y(0)
 \quad(t<T_*).
\]
The homogeneous Sobolev inequality now bounds
 $\|u(t)\|_6\leq C Y(0)^{1/2}$ uniformly.

For completeness, the continuation step is quantitative. In the mild
 equation, the heat--Leray estimate
\[
 \|e^{\nu\tau\Delta}\mathbb P\nabla\cdot F\|_6
 \leq C(\nu\tau)^{-3/4}\|F\|_3
\]
comes from heat-kernel differentiation, $L^3$--$L^6$ smoothing, and
 boundedness of the Leray projector on $L^3$. Since
 $\|u\otimes v\|_3\leq\|u\|_6\|v\|_6$, the mild bilinear term on a
 time interval of length $\delta$ has norm at most
 $C\nu^{-3/4}\delta^{1/4}\|u\|_{C_tL^6_x}\|v\|_{C_tL^6_x}$.
 The contraction argument therefore gives a lifespan bounded below by
 $c\nu^3M^{-4}$ whenever the initial $L^6$ norm is at most $M>0$.
 This is the local theory used in [E1]. Restarting near any proposed finite
 $T_*$ with the uniform bound just obtained extends the solution past it.
 Uniqueness of the strong solution makes the overlapping constructions agree.
 Smoothing for positive time, together with the initial Schwartz regularity,
 supplies the requested smoothness, including time zero. The energy identity
 supplies the global energy bound.\hfill$\square$

This subcase has the right invariance but is not the full target. Under
\[
 u_\lambda(x,t)=\lambda u(\lambda x,\lambda^2t),\qquad
 p_\lambda(x,t)=\lambda^2p(\lambda x,\lambda^2t),
\]
viscosity is unchanged and
 $E_\lambda=\lambda^{-1}E$, $Y_\lambda=\lambda Y$ at corresponding times.
 Thus $EY/\nu^4$ is invariant. A rescaling can make energy arbitrarily
 small while making enstrophy large; it cannot move arbitrary initial data
 into Proposition 5.1. Decreasing amplitude alone changes the nonlinear
 balance and creates different initial data.

\paragraph{A concrete conditional continuation target.}
Let $s_+(x,t)=\max\{0,\lambda_{\max}(S(x,t))\}$ and
 $b(t)=\|s_+(\cdot,t)\|_\infty$. Then $N\leq b(t)Y$, and
\[
 Y(t)\leq Y(0)\exp\left(2\int_0^t b(s)\,ds\right).
\]
Consequently finite accumulated positive strain up to a finite endpoint
 rules out blowup there by the same $L^6$ continuation argument.
 This criterion is proved here, but the integral is not controlled by energy.

A more selective estimate would suffice. Given $0<\theta<1$, suppose one
 can choose a nonnegative $k\in L^1(0,T_*)$ such that
\begin{equation}\label{ns:tail}
 \int_{\{s_+>k(t)\}}s_+(x,t)|\omega(x,t)|^2\,dx
 \leq\theta\nu Z(t)
 \quad\text{for almost every }t<T_*.
\end{equation}
The complementary region contributes at most $k(t)Y(t)$, giving
 $\tfrac12Y'+(1-\theta)\nu Z\leq kY$ and hence continuation.
 Unlike a global $L^\infty$ strain estimate, \eqref{ns:tail} allows large
 strain where its vorticity-weighted contribution can be absorbed by
 diffusion. Proving this for arbitrary initial data is an explicit missing
 geometric estimate; it is an additional assumption, not a consequence of
 incompressibility or a renamed energy bound.

\paragraph{Counterexample route: can the force be removed?}
Take a smooth forced reference flow $U$ and an unforced solution $v$ on a
 common regular interval, with equal initial data. Put $w=v-U$ and let $f$
 be the force of $U$. Subtraction gives
\[
 \partial_tw+(U\cdot\nabla)w+(w\cdot\nabla)U+(w\cdot\nabla)w
 =\nu\Delta w-\nabla q-f.
\]
The two transport terms vanish in the $L^2$ product with $w$, so
\[
 \tfrac12\frac d{dt}\|w\|_2^2+\nu\|\nabla w\|_2^2
 \leq\|\nabla U\|_\infty\|w\|_2^2+\|f\|_2\|w\|_2.
\]
Regularizing the norm at zero and applying Gronwall proves
\begin{equation}\label{ns:force}
 \|w(t)\|_2\leq
 \int_0^t\|f(s)\|_2
 \exp\left(\int_s^t\|\nabla U(r)\|_\infty\,dr\right)ds.
\end{equation}
The estimate gives no uniform control near a singular reference trajectory
 when the accumulated gradient diverges. Even a small $L^2$ difference
 would not preserve unbounded pointwise velocity, because concentrations
 can have arbitrarily small $L^2$ norm. Also, setting $f$ equal to a residual
 of a chosen ansatz must establish the prescribed force regularity; it does
 not make that residual zero. Thus deleting the force in [S4] is not a
 justified construction of a counterexample to (A).

\paragraph{Verification, remaining work, and outcome.}
The signs in the curl identity, the powers of $\nu$ in Young's inequality
 and the local lifespan, and the scaling of $E,Y$ were checked directly.
 The zero-datum and $Z=0$ cases were separated before division. The scalar
 blowup example satisfies only the inequalities claimed for it. No finite
 numerical simulation or check of a few Fourier modes could establish the
 all-data, all-time conclusion, and none is represented as doing so.

The strongest proved conclusion is the stated scale-invariant small-data
 subcase, together with the positive-strain and strain-tail continuation
 implications. For general Schwartz data the first missing step in this
 proof route is an a priori bound such as \eqref{ns:tail}, or another estimate
 controlling a critical continuation norm before the maximal time.
 The energy identity and \eqref{ns:stretch} cannot supply it by themselves.

Three concrete next targets are a vorticity-alignment estimate implying
 \eqref{ns:tail}; a bound on $L^\infty_tL^3_x$ permitting use of [E2]; and a
 stability estimate for a singular construction in a norm strong enough to
 transfer blowup while removing the force. Any candidate counterexample
 must satisfy the actual unforced equation, have Schwartz initial data and
 finite energy, and evade the continuation criteria above.

\textbf{UNRESOLVED.} No all-data regularity theorem or unforced singular
 solution is obtained. Written by GPT-6 Astra Ultra from the exact catalogue
 target, with primary-source scope checks and direct derivations; the
 argument is an independently written research attempt, not a verification
 of the announced forced proof.

\subsection{Sources}
\begin{itemize}
\item[S1]Charles L. Fefferman, Existence and Smoothness of the Navier-Stokes Equation, official Clay problem description.
\url{https://www.claymath.org/wp-content/uploads/2022/06/navierstokes.pdf}.
\textit{Formulation source. Consulted 24 September 2026: Statement (A), page 2.}
\item[S2]On the Navier–Stokes Millennium Prize Problem.
\url{https://openai.com/index/navier-stokes-solution/}.
\textit{Further reference; not a proof endorsement.}
\item[S3]Navier-Stokes Announcement — Clay Mathematics Institute.
\url{https://www.claymath.org/news/navier-stokes-announcement/}.
\textit{Further reference; not a proof endorsement.}
\item[S4]On the Navier–Stokes Millennium Prize Problem — analytical paper.
\url{https://cdn.openai.com/pdf/32d9f210-8b73-45e0-91bc-82a30aef8a9a/navier-stokes.pdf}.
\textit{Further reference; not a proof endorsement.}
\item[E1]Tosio Kato, \emph{Strong $L^p$-solutions of the Navier--Stokes equation in $\mathbb R^m$, with applications to weak solutions}, Mathematische Zeitschrift \textbf{187} (1984), 471--480.
\url{https://doi.org/10.1007/BF01174182}.
\textit{Local strong-solution framework; publisher record checked 27 September 2026. The $L^6$ continuation estimate used here is displayed in the argument.}
\item[E2]L. Escauriaza, G. A. Seregin, and V. \v Sver\'ak, \emph{$L_{3,\infty}$-solutions of the Navier--Stokes equations and backward uniqueness}, Russian Mathematical Surveys \textbf{58} (2003), 211--250.
\url{https://www.mathnet.ru/eng/rm609}.
\textit{Endpoint regularity theorem; journal record and abstract checked 27 September 2026. Here $L_{3,\infty}$ denotes $L^\infty$ in time with values in spatial $L^3$, not weak spatial $L^3$.}
\item[E3]Terence Tao, \emph{Finite time blowup for an averaged three-dimensional Navier--Stokes equation}, Journal of the American Mathematical Society \textbf{29} (2016), 601--674.
\url{https://arxiv.org/abs/1402.0290}.
\textit{Primary source for the averaged-nonlinearity obstruction; consulted 27 September 2026. It concerns a modified equation and is not an unforced Navier--Stokes counterexample.}
\end{itemize}
\end{document}
